\documentclass[12]{informe_investigacion}
\usepackage[spanish]{babel}
\usepackage[utf8]{inputenc}
\usepackage[a4paper, total={15.5cm, 24.7cm}]{geometry}
\usepackage{graphicx}
\graphicspath{ {images/} }

\title{Informe de Investigación - ULaSalle}
\author{WardoPo}
\date{March 2021}

\begin{document}

\begin{titlepage}

    \includegraphics[width=0.4\textwidth]{imagenes/La Salle Logo.png}
    
        \begin{center}
            \vspace*{1cm}
            \textbf{ESCUELA PROFESIONAL DE INGENIERÍA DE SOFTWARE}
            
            \vspace{1.5cm}
            [TITULO DEL  INFORME DE INVESTIGACIÓN]
            
            \vspace{1.5cm}
            [Autor]
            
            \vspace{1.5cm}
            [Grado a obtener]
            
            \vfill
            
            AREQUIPA – PERÚ
            
            \vspace{0.5cm}
            [Año]
        \end{center}

\end{titlepage}

\tableofcontents

\begin{resumen}
El resumen posee un conjunto de elementos los cuales dan una visión clara del trabajo. El resumen debe contener lo siguiente: 

\begin{itemize}
    \item Delimitación de la Investigación.
    \item El propósito de la investigación. 
    \item Forma de lograrlo. 
    \item Criterios que justifican al estudio.
    \item Fundamentación teórica empleada en el estudio. 
    \item Metodología de investigación empleada, se hace mención al tipo de estudio, el diseño de investigación, la modalidad de investigación (si es necesario),  técnica empleada para la recolección de datos, tipo de instrumento de recolección usado, validez y confiabilidad.
    \item Muy breve referencia a los resultados. 
    \item Señalamiento de las conclusiones más significativas 
\end{itemize}

La forma de redacción es en pasado impersonal, por ejemplo: una vez que recabamos la información (esta forma es incorrecta), en vez de ello se debe decir: una vez que se recabó la información (estas la forma correcta). 
\end{resumen}

\begin{abstract}
    
\end{abstract}

\begin{palabras_clave}
De 3 a 5 palabras clave del trabajo de investigación. 
\end{palabras_clave}

\chapter{Problema de Investigación}
Es la identificación de la problemática que se trata de solucionar por medio de la investigación y, para la tesis, es en sí la elección del tema que servirá de base para elaborarla mediante una proposición concreta indicando el contexto preliminar del problema, definición del problema, antecedentes y/o estado del arte preliminar, objetivos, justificación/relevancia.
El problema deberá cumplir una serie de condiciones que de alguna forma justifiquen el esfuerzo necesario para resolverlo. Entre ellas: originalidad, trascendencia, actualidad, relevancia y la posibilidad de permitir el uso de lo aprendido a lo largo de la carrera.

\section{Descripción del problema}
Descripción del contexto y definición del problema. Se describe la situación problemática del contexto desde la perspectiva científica. Debe enunciarse referencias que sustenten la problemática

\section{Objetivos de la Investigación}
Es la definición de lo que se pretende cumplir con la tesis. Debe existir una estricta correspondencia entre los objetivos, el planteamiento del problema y las conclusiones. También demanda una redacción sencilla, concreta y que contemple las siguientes reglas:
\begin{itemize}
    \item Iniciar el objetivo con un verbo en infinitivo.
    \item Determinar primero el qué se quiere y después el para qué se hace.
    \item Limitar la redacción a frases sustantivas.
\end{itemize}

%% La autoindentación dentro de subsecciones no es posible por lo que estas se manejan como ambientes separados.
\begin{subseccion}{Objetivo General}
Descripción de la finalidad principal que persigue la investigación, el motivo que le dará vigencia. El objetivo general y las preguntas de investigación están íntimamente relacionados entre si, por lo que deben ser coherentes
\end{subseccion}

\begin{subseccion}{Objetivos Específicos}
Señalan las actividades que se deben cumplir para avanzar en la investigación y lo que se pretende lograr en cada una de las etapas de ella, por ende, la suma de los resultados de cada uno de los objetivos específicos permitirá alcanzar el propósito integral del objetivo general. Especifica los logros concatenados que se pretende conseguir.
Para la formulación de los objetivos considere lo siguiente:
\begin{itemize}
    \item Deben estar dirigidos a los elementos básicos del problema
    \item Deben ser medibles y observables
    \item Deben ser claros y precisos
    \item Su formulación debe involucrar resultados concretos
    \item El alcance de los objetivos debe estar dentro de las posibilidades del investigador
    \item Deben ser expresados en verbos en infinitivos
\end{itemize}
\end{subseccion}

\section{Preguntas de Investigación}
Una vez que se tiene bien claro el problema, se redactan las preguntas de investigación de acuerdo al problema que se analizará. La pregunta de investigación es uno de los primeros pasos metodológicos que un investigador debe llevar a cabo cuando emprende una investigación. La pregunta de investigación debe ser formulada de manera precisa y clara, de tal manera que no exista ambigüedad respecto al tipo de respuesta esperado.
Las preguntas de investigación son operaciones mentales que hace el investigador para reconocer los puntos que le interesa abordar en su investigación. Por lo tanto, cuando las preguntas están planteadas incorrectamente el razonamiento lógico no entiende cuál es la operación que debe realizar. Las preguntas de investigación deben contener las siguientes características:
\begin{enumerate}
    \item Ser concretas: es decir no dar cabida a la vaguedad. Vaguedad significa que no se entiende exactamente por qué cosa pregunta.
    \item Ser claras: es decir dejar evidente lo que se pregunta.
    \item Ser precisas: es decir puntuales y exactas en lo que preguntan.
    \item Estar completas, es decir sobre todo que contengan sujeto o predicado.
    \item Siempre deben contener un adverbio de pregunta
\end{enumerate}

La pregunta de investigación puede ser una afirmación o un interrogante acerca del fenómeno, en forma precisa y clara, de tal forma que de ésta se desprendan los métodos, procedimientos e instrumentos.

Considere que No todas las investigaciones tienen hipótesis; todo depende del grado de conocimiento sobre el problema que se investiga. Sólo necesitan hipótesis las investigaciones que ya han rebasado la fase exploratoria y se encuentran en fase confirmatoria o verificatoria. Las hipótesis, son justamente el objeto de la confirmación o verificación. Intentar forzar la presencia de hipótesis cuando el conocimiento sobre un problema o la propia naturaleza de dicho problema no lo consienten es uno de los errores más frecuentes que se comenten en la práctica.

\chapter{Protocolo de Revisión}
\section{Estrategia de Búsqueda}
\section{Criterios de Selección de estudios}
\section{Estrategia para la extracción de datos}

\chapter{Resultados del Estado del Arte}
\section{Aspecto/Tema 1}
\section{Aspecto/Tema 2}
\section{Aspecto/Tema 3}

\begin{conclusiones}
Es la parte donde se manifiesta lo más destacado encontrado durante su investigación. Es una parte muy importante de la tesis puesto que en ella se indican los hallazgos y, en consecuencia, la comprobación de los objetivos. Aquí se muestran las aportaciones a la disciplina de estudio.
Es recomendable que el alumno elabore sus conclusiones tratando de cubrir por lo menos algunos de los siguientes puntos, mismos que sólo se presentan como una guía o referencia para una mejor elaboración de esas conclusiones, pero no se describen en el texto. Estas sugerencias se enfocan hacia:
\begin{itemize}
    \item Resultados encontrados.
    \item Demostración realizada.
    \item Comprobación (refutación) de la hipótesis.
    \item Conclusión general.
    \item Conclusiones parciales (útiles al trabajo).
    \item Aportaciones a su tema (disciplina).
\end{itemize}

Algunas consideraciones a tener en cuenta son:
\begin{itemize}
    \item Evitar que las conclusiones sean un resumen de cada capítulo pues en realidad se trata de consecuencias y determinaciones que conllevan una verdad.
    \item Su redacción debe ser clara, concreta, directa y enfática.
    \item Es importante hacer conclusiones específicas por cada punto de interés, pero sin abusar de este recurso. Todo estará en función al tema, su importancia y lo relevante de lo encontrado.
    \item Llegar a una conclusión global, si es posible, en que se concentren los aspectos fundamentales de la investigación, procurando abarcar solamente lo sustantivo y básico del tema.
    \item Evitar el tono imperativo e impositivo tanto como el tímido y desobligado.
\end{itemize}

\cite{einstein}

\end{conclusiones}

\bibliographystyle{IEEEtran}
\bibliography{bibliografia}

\end{document}
